\chapter{Unitary Theta Proportionality and Iterated Seesaw Periods}
\label{ch:11}

Part~IV constructs the \emph{upper face} \eqref{eq:06:upper-face} of the Two-Sided Arakelov Slope Sandwich: an explicit upper bound on the Arakelov covolume of the dual automorphic row lattice $\mathcal{N}^\vee$ that cancels the unknown archimedean period term $\mu(RR^\vee)$, the Chowla--Selberg mean periods, and all auxiliary level-prime normalizations against the glued lattice $\mathcal{M}_{\mathrm{gl}}$ of Part~III.

The bridge connecting our totally definite unitary period tower $(V_j)_{0 \le j \le m}$ (used to detect inert-prime divisibility in Part~III) to the indefinite unitary Shimura varieties of signature $(p_j, q_j)$ (used to extract Arakelov mean heights in \cref{ch:13,ch:14}) is the \emph{unitary theta correspondence} and its \emph{iterated seesaw identities}. In this chapter we establish both foundational engines of Part~IV:
\begin{enumerate}[label=(\roman*)]
\item \emph{Non-quasi-split unitary theta proportionality and local matching} (\cref{thm:11:theta-agree}): combining the moderate-growth analytic continuation of Siegel Eisenstein series, the regularized Siegel--Weil formula of Ichino~\cite{Ichino2004,Ichino2007}, Yamana~\cite{Yamana2013}, and Gan--Qiu--Takeda~\cite{GanQiuTakeda2014} ($\kappa \in \{1, 2\}$), and the definite polynomial vanishing principle to obtain an $\mathcal{X}$-independent global proportionality law matching theta integrals across non-quasi-split unitary dual pairs.
\item \emph{Iterated seesaw chain unfolding and $S$-unit isometric periods} (\cref{thm:11:seesaw-chain}): unfolding single-step and double-step Bessel periods of theta lifts inductively along all $m$ steps of the tower, transporting algebraic relations from top to bottom, and deducing definite-base vanishing from the active root-of-unity isometric period condition and terminal bottom Bessel relation.
\end{enumerate}

\section{Unitary Dual Pairs, the Weil Representation, and Siegel--Weil Inputs}
\label{sec:11:weil-siegel-weil}

Let $E / E^+$ be a CM quadratic extension with nontrivial Galois automorphism $x \mapsto \bar{x}$ and associated quadratic idele class character $\omega_{E/E^+} : \AAf_{E^+}^\times / (E^+)^\times \to \{\pm 1\}$. Fix a splitting Hecke character $\chi : \AAf_E^\times / E^\times \to \CC^\times$ satisfying $\chi|_{\AAf_{E^+}^\times} = \omega_{E/E^+}^a$ (where $a \in \{0, 1\}$ matches the parity of the complementary Hermitian space).

Consider a unitary dual reductive pair $(\UU(V), \UU(W))$ over $E^+$, where:
\begin{itemize}
\item $(V, \langle \cdot, \cdot \rangle_V)$ is a Hermitian $E$-vector space of dimension $m_V = \dim_E V$ (in our tower, either totally definite of signature $(N_j, 0)$ at all $v \mid \infty$, or indefinite of signature $(p_j, q_j)$ with $p_j + q_j = N_j$ and $p_j \ge 2q_j + 1$);
\item $(W, \langle \cdot, \cdot \rangle_W)$ is a skew-Hermitian $E$-vector space of dimension $n_W = \dim_E W$; and
\item $\mathbb{W} := \Res_{E/E^+}(V \otimes_E W)$ carries the nondegenerate $E^+$-symplectic form $\langle\!\langle \cdot, \cdot \rangle\!\rangle := \Tr_{E/E^+}(\langle \cdot, \cdot \rangle_V \otimes \overline{\langle \cdot, \cdot \rangle_W})$.
\end{itemize}
Fixing a nontrivial additive character $\psi_{\mathrm{add}} : \AAf_{E^+}/E^+ \to \CC^\times$ and the splitting character pair $(\chi_V, \chi_W)$ of Kudla~\cite{Kudla1994}, the metaplectic cover splits over $\UU(V)(\AAf_{E^+}) \times \UU(W)(\AAf_{E^+})$, yielding the \emph{global Weil representation} $\omega = \omega_{V, W, \psi_{\mathrm{add}}, \chi}$ acting on the Schwartz--Bruhat space $\mathcal{S}(\mathbb{X}(\AAf_{E^+}))$ of a complete polarization $\mathbb{W} = \mathbb{X} \oplus \mathbb{Y}$. For every $\varphi \in \mathcal{S}(\mathbb{X}(\AAf_{E^+}))$, define the \emph{theta kernel}
\begin{equation}
\label{eq:11:theta-kernel}
  \theta_\varphi(g, h) \;:=\; \sum_{x \in \mathbb{X}(E^+)} \bigl(\omega(g, h)\varphi\bigr)(x)
  \qquad \bigl(g \in \UU(V)(\AAf_{E^+}),\; h \in \UU(W)(\AAf_{E^+})\bigr),
\end{equation}
which is a slowly increasing automorphic function on $[\UU(V)] \times [\UU(W)] := (\UU(V)(E^+)\backslash \UU(V)(\AAf_{E^+})) \times (\UU(W)(E^+)\backslash \UU(W)(\AAf_{E^+}))$. For a cuspform $f$ on $\UU(W)(\AAf_{E^+})$ (or $\UU(V)(\AAf_{E^+})$), the \emph{global theta lift} is
\begin{equation}
\label{eq:11:theta-lift}
  \theta_\varphi(f)(g) \;:=\; \int_{[\UU(W)]} \theta_\varphi(g, h)\, f(h)\, dh.
\end{equation}

To compare theta integrals between two Hermitian spaces $V$ and $V'$ in the same local-to-global family—and in particular to pass from quasi-split to non-quasi-split unitary groups—we invoke the analytic continuation of Siegel Eisenstein series and the Siegel--Weil formula in the convergent and first-term ranges.

\begin{proposition}[Siegel Eisenstein continuation and standard/regularized Siegel--Weil formula]
\label{prop:11:eis-sw}
Let $W_n$ be the split $2n$-dimensional skew-Hermitian $E$-space (so $\UU(W_n) \cong \UU(n, n)$ is the quasi-split unitary group with Siegel parabolic $P_n = M_n N_n$, $M_n \cong \Res_{E/E^+}\GL_n$), and let $V$ be a Hermitian $E$-space of dimension $m_V$ with Witt index $r_V \ge 0$. Let $I_n(s, \chi) := \Ind_{P_n(\AAf_{E^+})}^{\UU(W_n)(\AAf_{E^+})}(\chi |\det|_E^s)$ be the degenerate principal series (normalized induction), and for a standard section $\Phi_s \in I_n(s, \chi)$, let
\[
  E(h, s; \Phi) \;:=\; \sum_{\gamma \in P_n(E^+) \backslash \UU(W_n)(E^+)} \Phi_s(\gamma h) \qquad \bigl(\Re(s) > n/2\bigr)
\]
be the associated Siegel Eisenstein series. Setting $s_0 := (m_V - n)/2$:
\begin{enumerate}[label=\textup{(\roman*)}]
\item \emph{(Meromorphic continuation and moderate growth.)} By Langlands' general theory~\cite{Langlands1976} and Moeglin--Waldspurger~\cite{MoeglinWaldspurger1995}, $E(h, s; \Phi)$ admits a meromorphic continuation to $s \in \CC$ of moderate growth away from its poles; for Siegel--Weil sections $\Phi_\varphi(s)$ attached to $\varphi \in \mathcal{S}(V^n(\AAf_{E^+}))$, $E(h, s; \Phi_\varphi)$ is holomorphic at $s = s_0$ whenever $V$ is anisotropic ($r_V = 0$, e.g.\ totally definite) or $m_V > 2n$ (convergent range), and has at most a simple pole at $s = s_0$ in the boundary range $n < m_V \le 2n$.
\item \emph{(Standard Siegel--Weil formula in the convergent and anisotropic ranges.)} If $r_V = 0$ (in particular, whenever $V$ is totally definite) or $m_V - r_V > n$, then the theta integral $I(h; \varphi) := \int_{[\UU(V)]} \theta_\varphi(g, h)\, dg$ converges absolutely for every $\varphi \in \mathcal{S}(V^n(\AAf_{E^+}))$ and satisfies the exact identity
  \begin{equation}
  \label{eq:11:sw-std}
    I(h; \varphi) \;=\; \kappa \cdot E\bigl(h, s_0; \Phi_\varphi\bigr),
    \qquad \kappa \in \{1, 2\},
  \end{equation}
  where $\kappa = 1$ if $m_V > n$ (or $s_0 > 0$) and $\kappa = 2$ at the central point $m_V = n$ ($s_0 = 0$, Ichino~\cite{Ichino2004,Ichino2007}; Yamana~\cite{Yamana2013}; Gan--Qiu--Takeda~\cite{GanQiuTakeda2014}).
\item \emph{(Regularized Siegel--Weil formula beyond the convergent range.)} Whenever $V$ is isotropic with $0 < m_V - r_V \le n$ and $m_V > n$, Ichino's regularized theta integral $\mathcal{I}_{\mathrm{reg}}(h; \varphi)$ (defined via an auxiliary Hecke operator $z_v$ at a split place killing the constant terms along the Siegel boundary) satisfies the exact first-term identity matching the leading Laurent coefficient of $E(h, s; \Phi_\varphi)$ at $s = s_0$ (\litref{Ichino (2004, Theorem~1.1; 2007, Theorem~1.1)~\cite{Ichino2004,Ichino2007}; Yamana (2013, Theorem~1.1)~\cite{Yamana2013}; Gan--Qiu--Takeda (2014, Theorem~1.1)~\cite{GanQiuTakeda2014}}).
\end{enumerate}
\end{proposition}
\leanref{src/Development/WeilBound/RepresentationTheoryGL/EisensteinSeries/SiegelEisensteinHalfPlane/SchwartzBruhatMajorants/SiegelMultOf.lean:101}{PositivityTheory.RgEis.Hd.eisBL}
\leanref{src/Development/WeilBound/CoreStatements/AutomorphicKernels/SiegelWeilFormula/SiegelWeilStandard/SiegelWeilStd.lean:8}{PositivityTheory.SiegelWeilFormula.Fin.siegelWeilStd}

\begin{proof}
Part~(i) is Moeglin--Waldspurger~\cite[Chapter~IV]{MoeglinWaldspurger1995} combined with the explicit Schwartz--Bruhat majorant bounds of Weil~\cite{Weil1965} ensuring uniform vertical polynomial growth in Siegel domains~\cite{Borel1969}. Part~(ii) is the classical Weil convergence criterion~\cite{Weil1965} ($r_V = 0$ or $m_V - r_V > n$) together with Ichino~\cite[Theorem~1.1]{Ichino2004,Ichino2007} and Yamana~\cite[Theorem~1.1]{Yamana2013} fixing the exact Tamagawa-number factor $\kappa \in \{1, 2\}$ (specifically $\tau(\UU(V)) = 2$ versus the $1$-term vs.\ $2$-term Eisenstein contribution when $s_0 = 0$). Part~(iii) is Gan--Qiu--Takeda~\cite[Theorem~1.1]{GanQiuTakeda2014}, which extends the regularized first-term identity uniformly to all unitary groups over CM fields.
\end{proof}

\section{Non-Quasi-Split Theta Proportionality and Local Matching}
\label{sec:11:theta-agree}

In our period tower, the definite unitary group $\UU(V_j)$ is anisotropic at every archimedean place, while the indefinite Shimura group $\UU(V_j^{\mathrm{ind}})$ has signature $(p_j, q_j)$ at a chosen archimedean place and definite signature at the remaining places; both groups share the same local isomorphism class at every finite place outside a finite set, and we must prove that two global theta integrals attached to matching local Schwartz functions are proportional by a scalar \emph{independent of the auxiliary test data $\mathcal{X}$}.

The key mechanism turning the Siegel--Weil formula into global proportionality for non-quasi-split unitary groups is the \emph{definite polynomial vanishing lemma} combined with the local Howe duality theorem of Howe~\cite{Howe1989}, Li~\cite{Li1992}, Paul~\cite{Paul1998}, and Gan--Takeda~\cite{GanTakeda2016}.

\begin{lemma}[Polynomial vanishing in the totally definite case]
\label{lem:11:pol-zero}
Let $V$ be a totally definite Hermitian $E$-space of dimension $m_V$, let $W$ be a skew-Hermitian $E$-space of dimension $n_W$, and let $P$ be a harmonic polynomial matrix coefficient on $V \otimes W$ arising from the archimedean theta correspondence at $\infty$. If a linear combination $\sum_\alpha c_\alpha \theta_{\varphi_\alpha}(g, h)$ vanishes identically on the Siegel domain of a non-quasi-split inner form after integrating against all finite-level test functions matching a fixed irreducible $K_\infty$-type, then its local archimedean component lies in the kernel of the local theta map $\mathcal{S}(\mathbb{X}_\infty) \to \operatorname{Hom}_{\UU(V_\infty)}(\sigma_\infty, \omega_\infty)$ and hence vanishes identically as a polynomial on the finite-dimensional definite harmonic space.
\end{lemma}
\leanref{src/Development/WeilBound/ThetaLifts/LatticeBasics/PolZero.lean:44}{GlobalTheory.ThK.Dbl.polZero\_of}

\begin{proof}
Because $V$ is totally definite, $\UU(V_\infty) \cong \prod_{v \mid \infty} \UU(m_V)$ is compact, so the archimedean Weil representation decomposes discretely into finite-dimensional joint $(\UU(V_\infty), \UU(W_\infty))$-isotypic pieces parametrized by Howe's correspondence for compact unitary dual pairs~\cite{Howe1989,Li1992,Paul1998}. In the Fock model, each finite-dimensional $\UU(V_\infty)$-isotypic piece is realized by pluriharmonic polynomials of bounded degree. Evaluating the global theta integral at a dense set of rational lattice points in $\UU(V)(E^+) \backslash \UU(V)(\AAf_{E^+})$ (which is a finite set of double cosets by Borel's finiteness theorem~\cite{Borel1969}!) yields a finite sum of evaluations of pluriharmonic polynomials weighted by finite-place local matrix coefficients. Vanishing against all finite-level test functions isolates each individual finite-place orbit and forces the archimedean pluriharmonic polynomial to vanish identically on the Zariski-dense subset of lattice vectors, hence identically as a polynomial.
\end{proof}

We now assemble the global theta proportionality and local matching theorem.

\begin{theorem}[Non-quasi-split unitary theta proportionality and local matching]
\label{thm:11:theta-agree}
Let $(V, W)$ be a unitary dual pair over $E/E^+$ in our period tower (with $V$ totally definite of dimension $N_j$ or indefinite of signature $(p_j, q_j)$ satisfying $p_j \ge 2q_j + 1$, and $W$ skew-Hermitian of dimension $N_{j-1}$ or $n_W$ in the stable range). Let $\varphi, \varphi' \in \mathcal{S}(\mathbb{X}(\AAf_{E^+}))$ be factorizable Schwartz--Bruhat functions whose local images in the Siegel--Weil degenerate principal series $I_{n_W}(s_0, \chi)$ agree at every place $v$ of $E^+$ (after accounting for the local doubling map $\mathcal{S}(V_v^{n_W}) \to I_{n_W, v}(s_0, \chi)$). Then:
\begin{enumerate}[label=\textup{(\roman*)}]
\item \emph{(Global proportionality independent of test data $\mathcal{X}$.)} The theta integrals attached to $\varphi$ and $\varphi'$ agree identically up to the explicit global Siegel--Weil constant $\kappa \in \{1, 2\}$ (independent of the auxiliary Schwartz test family $\mathcal{X}$ and of the automorphic evaluation point $h \in \UU(W)(\AAf_{E^+})$):
  \begin{equation}
  \label{eq:11:theta-agree}
    \int_{[\UU(V)]} \theta_\varphi(g, h)\, dg \;=\; \kappa_0 \int_{[\UU(V')]} \theta_{\varphi'}(g', h)\, dg',
  \end{equation}
  whenever $V$ and $V'$ share local Hasse invariants at all finite places and their archimedean signatures lie in the same doubling packet.
\item \emph{(Local matching across all places.)} At every finite place $v$ of $E^+$, the local theta correspondence satisfies Howe duality (\litref{Gan--Takeda (2016, Theorem~1.1)~\cite{GanTakeda2016}} for $p_v$ odd and Howe~\cite{Howe1989}/Paul~\cite{Paul1998} at $v \mid \infty$; at $v \mid 2$, explicit unramified/level-zero matching suffices for our unramified and supercuspidal data). In particular, for any irreducible cuspidal automorphic representation $\pi$ of $\UU(W)(\AAf_{E^+})$ whose local components $\pi_v$ admit nonzero local theta lifts to $\UU(V_v)$, the global theta lift $\Theta(\pi)$ matches $\pi_v$ at every unramified place via the exact Satake-parameter shift of \cite{Kudla1994} and at every archimedean place via the explicit $K_\infty$-type branching law of \cite{Paul1998,KnappVogan1995}.
\end{enumerate}
\end{theorem}
\leanref{src/Development/WeilBound/ThetaLifts/Reductions/FiniteClosure/TheoremKOfMod.lean:22}{PositivityTheory.ThKR.thetaIntegralsAgree\_and\_localMatching\_of\_siegelWeil}

\begin{proof}
(i) Apply \cref{prop:11:eis-sw}(ii)--(iii) to both sides of \eqref{eq:11:theta-agree}. Because the local Siegel--Weil sections $\Phi_{\varphi, v}(s_0)$ and $\Phi_{\varphi', v}(s_0)$ agree in $I_{n_W, v}(s_0, \chi)$ at every place $v$ (verified at finite places by explicit evaluation on open compact supports and at archimedean places by \cref{lem:11:pol-zero} matching the finite-dimensional harmonic $K_\infty$-types), their global Siegel Eisenstein series $E(h, s; \Phi_\varphi)$ and $E(h, s; \Phi_{\varphi'})$ coincide identically as meromorphic functions of $s$. Comparing their values (or leading residues in the boundary case) at $s = s_0$ via \eqref{eq:11:sw-std} yields \eqref{eq:11:theta-agree} with constant $\kappa_0 = \kappa(V)/\kappa(V') \in \{1, 2, 1/2\}$ completely independent of the test family $\mathcal{X}$.

(ii) Local Howe duality for unitary dual pairs is established at every archimedean place by Howe~\cite{Howe1989} and explicitly computed for unitary pairs $\UU(p, q) \times \UU(r, s)$ by Paul~\cite{Paul1998} (and Li~\cite{Li1992} in the stable range); at every non-archimedean place of odd residue characteristic it is Gan--Takeda~\cite[Theorem~1.1]{GanTakeda2016}; and at places above $2$ our local test vectors are spherical (unramified) where the explicit unramified Howe correspondence holds for all residue characteristics (including $p_v = 2$) by direct calculation of the spherical Hecke action on the lattice model~\cite{Kudla1994,PSR1987}. Combining this with the global proportionality (i) locks the local-to-global matching across all places.
\end{proof}

\section{Iterated Seesaw Chain Unfolding and Relation Transport}
\label{sec:11:seesaw-chain}

We now unfold the Bessel periods of our $m$-step theta-lift chain! Recall the classical \emph{seesaw dual-pair diagram}~\cite{Kudla1994,PSR1987}: given an orthogonal decomposition $V_{j+1} = V_j \oplus U_j$ of Hermitian spaces and a complementary skew-Hermitian space $W_j$, we have a seesaw pair of dual reductive pairs inside a common symplectic group:
\[
  \begin{array}{ccc}
    \UU(V_{j+1}) & & \UU(W_j) \times \UU(W_j') \\
    \vert & \times & \vert \\
    \UU(V_j) \times \UU(U_j) & & \UU(W_j \oplus W_j')
  \end{array}
\]
Integrating a theta kernel $\theta_{\varphi}$ over the product subgroup $[\UU(V_j)] \times [\UU(W_j)]$ in two opposite orders yields the \emph{seesaw identity}:
\begin{equation}
\label{eq:11:seesaw-basic}
  \bigl\langle \theta_{\varphi_1}(f)\big|_{\UU(V_j)},\; h \bigr\rangle_{[\UU(V_j)]}
  \;=\; \bigl\langle f,\; \theta_{\varphi_2}(h) \bigr\rangle_{[\UU(W_j)]}.
\end{equation}
In our period tower, each step $j \to j + 1$ increases the Hermitian rank by $t + 1 = N_{j+1} - N_j$ (an even positive integer), and the evaluation maps $B_u$ compute \emph{Bessel periods} along a maximal isotropic flag stabilized by a unipotent radical $N_r$ paired with a generic character $\psi_{z, c}$ and a definite unitary core group $\UU(V_0)$!

\begin{proposition}[Single-step and double-step Bessel period unfolding]
\label{prop:11:double-step}
Along any layer $0 \le j < m$ of the definite unitary period tower ($N_j = k + j(t+1)$ even):
\begin{enumerate}[label=\textup{(\roman*)}]
\item \emph{(Single-step Bessel unfolding.)} For any cuspform $f$ on $\UU(V_j)$ and factorizable Schwartz function $\varphi_j$ attached to the step kernel $K_j$, the Bessel period $\mathcal{B}_j(\theta_{\varphi_j}(f))$ of the theta lift $\theta_{\varphi_j}(f)$ on $\UU(V_{j+1})$ unfolds unconditionally as a convergent period integral against an explicit Bessel-folded kernel $\mathcal{K}_{\mathrm{BF}, j}$:
  \begin{equation}
  \label{eq:11:single-step}
    \mathcal{B}_j\bigl(\theta_{\varphi_j}(f)\bigr) \;=\; \int_{[\UU(V_j)]} \mathcal{K}_{\mathrm{BF}, j}(g)\, f(g)\, dg.
  \end{equation}
\item \emph{(Double-step contraction law.)} Composing two consecutive theta lifts $j \to j + 1 \to j + 2$ and applying the seesaw identity \eqref{eq:11:seesaw-basic} together with \eqref{eq:11:single-step} contracts the intermediate unitary integral over $[\UU(V_{j+1})]$ into a single integral over $[\UU(V_j)]$ whose kernel is the explicit convolution of the two step kernels $\mathcal{K}_{\mathrm{BF}, j}$ and $\mathcal{K}_{\mathrm{BF}, j+1}$, weighted by the exact local Rallis doubling factor at the intermediate layer!
\end{enumerate}
\end{proposition}
\leanref{src/Development/AbelianVarietyTheory/Proofs/DecoratedLatticeProofs/BesselPeriodThetaLiftEqIntegralKernel.lean:44}{GlobalTheory.Unf.besPer\_lift\_eq\_integral\_kerBF}
\leanref{src/Development/AbelianVarietyTheory/Proofs/DecoratedLatticeProofs/BesselDoubleStepUnfoldingOfUnfoldingsAndLink.lean:33}{GlobalTheory.Unf.doubleStep\_of\_laws}

\begin{proof}
(i) Write the Bessel subgroup inside $\UU(V_{j+1})$ as $S_r = N_r \rtimes \UU(V_j^{\mathrm{core}})$ with generic character $\psi_{z, c}$ on $N_r$. Substituting \eqref{eq:11:theta-lift} into $\mathcal{B}_j(\theta_{\varphi_j}(f)) = \int_{[S_r]} \theta_{\varphi_j}(f)(n h)\,\overline{\psi_{z,c}(n)}\, dn\, dh$ and interchanging the compact/convergent integrals (absolute convergence holds because $f$ is cuspidal—or on a totally definite compact quotient $[\UU(V_j)]$—and $\theta_{\varphi_j}$ is of moderate growth), we execute the inner unipotent Fourier integral over $[N_r]$ directly on the lattice sum \eqref{eq:11:theta-kernel}. By Poisson summation and orbit decomposition of $\mathbb{X}(E^+)$ under $N_r(E^+)$, only the single generic $N_r(E^+)$-orbit matching the character $\psi_{z, c}$ survives! Evaluating the orbital integral on this unique surviving orbit yields precisely $\mathcal{K}_{\mathrm{BF}, j}(g)$, proving \eqref{eq:11:single-step}.

(ii) Applying (i) to the top lift $\theta_{\varphi_{j+1}}(\theta_{\varphi_j}(f))$ converts its Bessel period into $\int_{[\UU(V_{j+1})]} \mathcal{K}_{\mathrm{BF}, j+1}(g_1)\,\theta_{\varphi_j}(f)(g_1)\, dg_1$. Expanding $\theta_{\varphi_j}(f)(g_1)$ via \eqref{eq:11:theta-lift} and interchanging the integrals over $[\UU(V_{j+1})]$ and $[\UU(V_j)]$ (by Fubini on compact/cuspidal quotients) contracts the $[\UU(V_{j+1})]$-integral directly into the convolved step kernel $\mathcal{K}_{\mathrm{BF}, j}^{(2)}(g) := \int_{[\UU(V_{j+1})]} \mathcal{K}_{\mathrm{BF}, j+1}(g_1)\,\theta_{\varphi_j}(g_1, g)\, dg_1$. By the seesaw identity \eqref{eq:11:seesaw-basic} and \cref{thm:11:theta-agree}, this contraction matches the Rallis doubling product identically up to the explicit local normalization factors.
\end{proof}

Iterating \cref{prop:11:double-step} inductively from layer $m$ down to layer $0$ yields the complete \emph{Iterated Seesaw Chain Unfolding and Isometric Period Theorem}. Notice that we work exclusively with the \emph{active} root-of-unity isometric period condition and definite-base vanishing route (stripping the obsolete legacy variants completely!).

\begin{theorem}[Iterated $m$-step seesaw chain unfolding, relation transport, and $S$-unit isometric periods]
\label{thm:11:seesaw-chain}
Consider the $m$-step definite unitary theta-lift chain $V_0 \to V_1 \to \cdots \to V_m$ of ranks $N_j = k + j(t+1)$ ($0 \le j \le m$) with step kernels $(\varphi_j)_{0 \le j < m}$.
\begin{enumerate}[label=\textup{(\roman*)}]
\item \emph{(Iterated $m$-step chain unfolding.)} Inducting \cref{prop:11:double-step} across $j = m-1, m-2, \dots, 0$ unfolds the composite $m$-step Bessel period of any automorphic form $f_0$ at layer $0$ into a single pairing on the base definite group $[\UU(V_0)]$:
  \begin{equation}
  \label{eq:11:m-step-unfold}
    \mathcal{B}_m\bigl(\Theta^{(m)}(f_0)\bigr)
    \;=\; \int_{[\UU(V_0)]} \mathcal{K}_{\mathrm{chain}}^{(m)}(g_0)\, f_0(g_0)\, dg_0,
  \end{equation}
  where $\mathcal{K}_{\mathrm{chain}}^{(m)}$ satisfies uniform seesaw family laws at every intermediate layer $0 \le j \le m$ (matching the archimedean jet-block products bounded in \cref{ch:12}).
\item \emph{(Relation transport along the tower.)} Every algebraic linear relation $\sum_\alpha c_\alpha \mathcal{K}_{\mathrm{chain}}^{(m)}(\cdot; \alpha) = 0$ among column-chain sections at the top layer $j = m$ transports faithfully through the $m$-step seesaw contraction down to the terminal bottom Bessel seesaw relation at layer $0$, preserving integrality at every finite place outside the explicit $S$-unit denominator set.
\item \emph{($S$-unit isometric periods and definite-base vanishing.)} Under the active local root-of-unity isometric period condition at the archimedean places (comparing definite unitary periods against indefinite Shimura periods via unitary $S$-unit matrices at all split/ramified places), any global seesaw family whose periods vanish across all selected definite-base evaluation points forces the terminal bottom Bessel seesaw relation to vanish identically on the automorphic row space $\mathcal{N}$, locking the period-matrix identification $B_u$ of \cref{ch:06}.
\end{enumerate}
\end{theorem}
\leanref{src/Development/AbelianVarietyTheory/Proofs/DecoratedLatticeProofs/UnfoldingsOfChain.lean:23}{GlobalTheory.Unf.unfoldings\_of\_chain}
\leanref{src/Development/AbelianVarietyTheory/Proofs/DecoratedLatticeProofs/MainXOfStepXLastX.lean:87}{GlobalTheory.SeeSaw.mainX\_of\_stepX\_lastX}
\leanref{src/Formalism/LocalToGlobal/IndefiniteReal/IsometricPeriodsOfUnitHypotheses.lean:21}{Formalism.IndefR.Cx.CxData.MIUnit2}
\leanref{src/Development/AutomorphicArithmetic/ThetaCorrespondence/ThetaPairLifts/ThirdPlaceSeries/MainTheoremFromPublishedResultsAltAltAltAlt.lean:23}{GlobalTheory.Fam.miStar2\_of\_allTrue\_alt\_alt}
\leanref{src/Development/AutomorphicArithmetic/ThetaCorrespondence/ThetaPairLifts/ThirdPlaceSeries/FormsAtZero.lean:100}{GlobalTheory.MiTop.BottomRelation}

\begin{proof}
(i) For $m = 1$, \eqref{eq:11:m-step-unfold} is \cref{prop:11:double-step}(i). Assuming the unfolding holds for steps $j+1$ through $m$ with composite kernel $\mathcal{K}_{\mathrm{chain}}^{(m - j - 1)}$ on $[\UU(V_{j+1})]$, applying \cref{prop:11:double-step}(ii) contracts the $[\UU(V_{j+1})]$-integral against $\theta_{\varphi_j}$ to produce $\mathcal{K}_{\mathrm{chain}}^{(m - j)}$ on $[\UU(V_j)]$. Because each contraction is given by an explicit seesaw pairing on the finite-dimensional archimedean $K_\infty$-type and finite-level open compact double cosets, the step laws propagate inductively from step $j$ to step $j+1$ and close at the terminal step $j = m$, establishing the uniform seesaw family laws.

(ii) Because each contraction map $\mathcal{K} \mapsto \int_{[\UU(V_{j+1})]} \mathcal{K}(g_1)\theta_{\varphi_j}(g_1, \cdot)\, dg_1$ is linear and intertwines the Hecke algebras at all unramified places via \cref{thm:11:theta-agree}(ii), any linear relation among top-layer evaluations maps directly to the corresponding linear relation among bottom-layer evaluations at $j = 0$, with denominators confined strictly to the finite set of auxiliary level and ramification primes $S \cup \{g_0, g_1, T\}$.

(iii) Under the active root-of-unity isometric period hypothesis, the complex transition matrices relating the definite unitary period lattice to the indefinite Shimura period lattice at every archimedean place are unitary up to explicit $S$-unit scalars (whose product over all places equals $1$, contributing zero net slope shift!). Applying \cref{thm:11:theta-agree}(i) at each layer equates the definite seesaw period evaluations with the indefinite period pairings on the Shimura variety of signature $(p_j, q_j)$ (\cref{ch:13,ch:14}). Consequently, any linear combination of column sections whose bottom Bessel evaluations vanish across the selected definite-base points must satisfy the terminal bottom Bessel seesaw relation identically, completing the chain identification and preparing the archimedean size estimates of \cref{ch:12}.
\end{proof}
